% 2010 年北京卷理科第 5 题：正视图、侧视图与四个俯视图选项。
% 来源 PDF：普通高考/2010/2010北京理.pdf，SHA-256 6b5819672aaf59aecfce92e08c2d045a80575803527d61813497dd669a8d5529。
% 原页：第 1 页（A3 三栏）第 1 栏中部，书页码 575。
\begin{tikzpicture}[semithick, every node/.style={inner sep=1.2pt, font=\small}]
  % 正视图
  \draw (0,0) rectangle (1.6,1.0);
  \draw (0,0.6) -- (0.6,0.6) -- (0.6,1.0);
  \node[below] at (0.8,-0.1) {正（主）视图};
  % 侧视图
  \draw (2.4,0) rectangle (4.0,1.0);
  \draw (2.4,0.6) -- (3.4,0.6) -- (3.4,1.0);
  \node[below] at (3.2,-0.1) {侧（左）视图};
  % 四个俯视图选项
  \draw (0,-2.0) rectangle (1.6,-0.8);
  \draw (0.6,-2.0) -- (0.6,-1.4) -- (1.6,-1.4);
  \node[below] at (0.8,-2.1) {（A）};
  \draw (2.4,-2.0) rectangle (4.0,-0.8);
  \draw (2.4,-1.4) -- (3.4,-1.4) -- (3.4,-0.8);
  \node[below] at (3.2,-2.1) {（B）};
  \draw (0,-3.6) rectangle (1.6,-2.4);
  \draw (0,-3.0) -- (0.6,-3.0) -- (0.6,-2.4);
  \node[below] at (0.8,-3.7) {（C）};
  \draw (2.4,-3.6) rectangle (4.0,-2.4);
  \draw (3.4,-3.6) -- (3.4,-3.0) -- (4.0,-3.0);
  \node[below] at (3.2,-3.7) {（D）};
\end{tikzpicture}
